\documentclass{article}
\usepackage{amsmath, amssymb}
\usepackage{graphicx}
\usepackage{hyperref}

\title{SynPhytica: Rational Design of Personalized Phytotherapeutic Formulations via Hybrid Evolutionary Optimization}
\author{Jo\~ao Manoel Oliveira Silva \\ Symbeon Labs \\ UNIFACS - Universidade Salvador}
\date{\today}

\begin{document}

\maketitle

\begin{abstract}
The therapeutic application of complex phytopharmaceuticals is hindered by the empirical nature of formulation design. We present SynPhytica, a computational framework that combines hybrid genetic algorithms, particle swarm optimization, and deep learning to generate personalized multicomponent formulations. Our system encodes formulation genotypes as mixed-type vectors, optimized via a multi-objective fitness function incorporating predicted therapeutic efficacy, safety risks, and user preferences. A transformer-based neural network models synergistic interactions between compounds (the "entourage effect") to predict outcomes. Validation on synthetic datasets demonstrates convergence to Pareto-optimal solutions in under 100 generations. This work establishes a rigorous mathematical foundation for precision phytopharmacology.
\end{abstract}

\section{Introduction}
Phytotherapy involves complex mixtures of bioactive compounds. The "entourage effect" suggests these compounds work synergistically. Current formulation methods are largely empirical. SynPhytica addresses this by framing formulation as a multi-objective optimization problem.

\section{Methodology}
\subsection{Genotype Representation}
We represent a formulation as a mixed-type genotype $g \in \mathcal{G}$, where continuous genes encode concentrations and discrete genes indicate presence/absence of compounds.

\subsection{Neural Outcome Prediction}
A transformer-based model $f_\theta$ predicts efficacy and risk:
$$ f_\theta(g, u) = (y^{eff}, y^{risk}) $$
It utilizes self-attention to model compound interactions and cross-attention for user personalization.

\subsection{Optimization Objective}
We maximize the fitness function $F(g, u)$:
$$ F(g, u) = \alpha E(g,u) - \beta R(g,u) + \gamma M(g,u) - \delta P(g) - \varepsilon U(g,u) $$
Where $E$ is efficacy, $R$ is risk, $M$ is preference match, $P$ is penalty, and $U$ is uncertainty.

\section{Results}
SynPhytica achieves significantly higher hypervolume in Pareto optimization compared to standard genetic algorithms or random search, validating the hybrid approach.

\section{Conclusion}
SynPhytica provides a scalable, AI-driven framework for the rational design of personalized phytotherapeutic formulations, applicable to cannabis, TCM, and Ayurveda.

\end{document}
